\vspace{-0.5cm}
\section{Challenges} \label{sec:challenges}
\vspace{-0.2cm}

% \todo{say challenges span many of these tasks; (Make the transition from sec 2 to sec 3 smooth)}

While the field of AI for code has made fruitful progress, cutting-edge AI still struggles with SWE tasks, especially at larger scopes and higher levels of algorithmic complexity. Next, we discuss ten key challenges in AI for code. Each challenge spans multiple tasks, and progress on any can lead to improvements on many tasks at once. We order these roughly by how easy it is to resolve them.

% A major bottleneck in the adoption of LLM systems for real-world applications is whether they are performant enough to outweigh the pros (time and labor) against the cons (codebase/developer burden). This is challenged by multiple aspects:

\vspace{-0.2cm}

\subsection{Evaluation and Benchmarks}

\vspace{-0.1cm}

Our taxonomy of tasks and measures highlights some of the shortcomings of today's evaluations and benchmarks. For example, the majority of today's coding evaluations have no level of human intervention, with a few, such as Copilot-Arena \citep{chi2024copilot}, having low to medium autonomy. 
% \ms{nit: copilot-arena style evals}
HumanEval, MBPP, APPS, CodeContests, and LiveCodeBench are all at function-level scope, with low to medium-high algorithmic complexity. Commit0 and SWE-Bench, TestGenEval~\citep{jain2024testgeneval}, RefactorBench~\citep{gautam2024refactorbench} are at project-level scope with low to medium algorithmic complexity. Today, we lack benchmarks 1) at project-level scope with high algorithmic complexity, 2) for many tasks other than code generation, and 3) requiring human intervention. Next, we note that current coding evaluations primarily focus on the code generation task. 
Most of the tasks discussed in Sec~\ref{sec:challenges} are either not studied such as Code QA or only studied in self-contained scopes like EvalPerf~\citep{evalperf}, formal verification~\citep{sun2024clover}.

% \textbf{Construct Validity}:
% Construct validity refers to how closely a measurement reflects the underlying concept. 
% Given the implications of rapid performance improvement in the AI for the code domain, it is essential to have high-construct validity benchmarks evaluating how well programming agents can perform.
% While benchmarks like SWE-Bench come close, user experiences do not currently match rapid performance gains obtained from them. 

% Existing works that focus on 

% \todo{talk about challenge in programing systems community to measure the success of new PLs. notorously dificult to eavluate. User studies fail to capture maintainability, learning effects.}



%\todo{I wonder if it makes sense to make a list of benchmarks and which category they belong to}

% Given the rapid pace of progress of LLMs, 

% Machine learning progress is in a large part guided by making steady progress (``hill climbing'') on some set of benchmarks.
% For example, \naman{ImageNet served ...}.
% Similarly, AI for code has witnessed a rising interest in the development of better coding evaluations

% Thus the choice of benchmarks developed and used in this process 

% Machine learning approaches are typically about optimization and progress is made through improving (``hill climbing'') on benchmarks. 
% Therefore, developing appropriate benchmarks that allow generalization to 
% the choice of benchmarks developed or used during the machine learning life-cycle is very important.
% Here, we draw upon learnings from prior works and present directions for 

% The rate of progress in artificial intelligence is considerably aff

\vspace{-0.2cm}

\subsection{Intelligent Tool Usage} \label{sec:tool-use}
\vspace{-0.1cm}

Software engineering has witnessed the development of various open and proprietary tooling support for programming, debugging, analysis, and code management throughout time.
% Software engineering is a long-standing field and has support for mature tooling support f
% which has led to development of various mature tooling to support developers in programming and debugging tasks. 
For example, linters and type-checkers provide assurances on static code correctness. 
Print statements and debuggers are used for dynamically analyzing and debugging programs.
Beyond programming, such tools are richly integrated into the entire software development lifecycle, e.g. code navigation or search, reviewing code, CI testing.
While many efforts combine LLMs and agents with tools, they do not achieve fully dynamic and intelligent software engineering tool usage. There are a few challenges towards this goal: first, the agent must identify which tools could potentially be useful for the task at hand. Second, the agent then needs to decide when to invoke the tool. Third, the agent then must figure out how to best make use the tool. Finally, the agent needs to incorporate the output provided by the tool in order to inform its next steps.


% There are a few challenges towards achieving dynamic and intelligent software engineering tool usage. First, the agent must identify which tools could potentially be useful for the task at hand. The agent then needs to decide when to invoke the tool. A complex debugging task might require the use of \texttt{pdb} or \texttt{gdb} to track intermediate program states, while looking at input-output pairs may be sufficient for simple debugging tasks. The agent then must figure out how to invoke the tool. If the agent knows that a certain function in a program has an error, it may wish to walk through only that function instead of the entire code from start to finish. Finally, the agent needs to incorporate the output provided by the tool in order to inform its next steps, e.g. edit the code if a bug was uncovered or run the tool again otherwise.

% \textit{Example: Performance Instrumentation}: CSI \citep{schardl2017csi} is a way to instrument software by programatically inserting hooks to track objects such as memory loads/stores and function entry/exits to measure code coverage and performance. For an LLM agent to use CSI effectively, it must 1) learn the CSI API, 2) decide when and how to use CSI (such as when it writes a suspected performance bottleneck), and 3) use the CSI output to decide how to proceed.


% \textit{Example: CSI}: A common way to instrument software systems is known as \textit{compiler-inserted program instrumentation}. CSI \citep{schardl2017csi} is a tool that inserts instrumentation hooks to track objects such as memory loads/stores, function entry/exits, and basic blocks. CSI contains tools like code coverage tools, a memory-operations counter, a performance profiler, and a call-graph generator. To use the tool, the user must follow the API in order to write hooks so the correct aspects can be profiled. Tools like CSI are very valuable when trying to improve the performance of a piece of code, but are not trivial to use. In order for an LLM agent to use CSI effectively, it must first familiarize itself with the CSI API. Then, it needs to know exactly which aspects of the code to instrument, such as placing hooks before and after a function suspected to be a bottleneck. Finally, the agent needs to learn how to use the output of the tool to inform its approach to the task, such as deciding whether a block of code can be further optimized after seeing its performance profile.




% There are a lot of programming and debugging tools that people use that most AI programming systems probably don't leverage or know how to use effectively:

% \begin{itemize}
%     \item Print statement debugging.
%     \item Debuggers (pdb/gdb).
%     \item Linters.
%     \item For optimization, things like Valgrind.
% \end{itemize}

\vspace{-0.2cm}

\subsection{Human-AI Collaboration} \label{sec:c-hai-colab}

\vspace{-0.1cm}

While AI coding assistants are increasingly more powerful, the majority of AI coding assistants are still at a low to medium autonomy level, with constant human supervision required. We identify a few key challenges of today's AI coding systems that prevent these systems from working with humans effectively at higher levels of autonomy.

\textbf{Lack of controllability}: When programmers use AI coding assistants, they are often looking for a specific result and functionality. However, they currently do not have reliable ways of steering LLMs to generate a very specified chunk of code. Often, they rely on a guess-and-check like approach where the LLM is repeatedly sampled or given feedback until outputting a piece of code the programmer likes. As a result, humans still need to spend a lot of time reviewing and modifying the code to ensure that it performs their desired functionality \citep{weisz2024examining}. 

\textbf{Identifying when human input is necessary}: LLMs rarely defer to humans for clarification, while developers often ask questions to clarify the description of a task provided by their peers. One example is a product manager refining a requirements document: developers confused about the requirements or the scope ask questions and add comments to the document, which are typically resolved by the manager with the goal of iteratively disambiguating the requirements \citep{nahar2022collaboration}. Based on its knowledge of existing software, AI should be able to incorporate implicit priors from a specification while keeping the user in the loop. For example, when designing an academic website for someone, there are implicit requirements, such as including a person's list of publications and contact information, but whether to include the person's GPA may be a clarification point.

% \textbf{Maximizing information from human interactions}: 
% - Identifying when human input is necessary

% - Maximizing information from human interactions (pragmatic stuff)

% - Support for multimodal (Zoom call with AI agent + screen share)

% {\color{blue}Fully autonomous software development remains an open challenge. Therefore, a better interface / workflow for human oversight is essential to integrate AI into development workflows effectively.}



% \begin{itemize}
%     % \item Lack of controllability, developers don't have reliable ways of model steering. Things are very guess and check right now, sample and test.
%     % \item From HCI perspective, we haven't found the right interface between humans and programmers.
%     % \item Generally, humans still have to check if AI copilot-written code is correct, people often say they spend more time debugging AI code than it would have taken to just write it themselves.
%     \item Still a lot of differences in human workflow vs. AI agent workflow. e.g. when debugging a pipeline with a large transformer, a human might use a smaller transformer to make loading time faster.
%     \item How can AI be used to augment humans rather than replace humans?
% \end{itemize}

\vspace{-0.2cm}

\subsection{Vague Specifications and User Misalignment} \label{sec:c-specifications}

\vspace{-0.1cm}

% \todo{discuss testing}

When using code LLMs, we typically prompt them with a natural language specification. This can include a natural language description of the desired code, input-output examples, relevant code snippets, and other functional requirements. However, there is a gap in the level of abstraction between English and code, leading to incomplete or ambiguous specifications. For longer programs, the number of ambiguous decision points also increases, and choices traditionally made by a human are now 
implicitly incorporated into the LLM's generated code. This often leads to user misalignment due to vague specifications.

\textbf{Inherent trade-offs in software development}: Designing large software systems always surfaces trade-offs between various desiderata such as readability, scalability, performance, maintainability, reliability, security, etc. These trade-offs are often context dependent. A long-term and rapidly moving project may be willing to trade off some performance to have simplicity and readability. Performance-critical applications may completely sacrifice readability to eke out every millisecond of speed. 
% (such as using \href{https://graphics.stanford.edu/~seander/bithacks.html}{bit-twiddling hacks}) 
Finding the sweet spot among these trade-offs can often involve extensive prototyping and benchmarking to understand the performance characteristics of different approaches. User specifications rarely include details about these trade-offs, nor do models often take them into account.

% \textbf{Implicit constraints}: Aside from functional/semantic correctness, there are also often implicit constraints in writing code. For example, many companies such as \href{https://opensource.janestreet.com/standards/}{Jane Street} and \href{https://google.github.io/styleguide/}{Google} have style guides, and many GitHub repositories explicitly outline style elements that new code ought to follow. \cite{zou2019does} find that GitHub pull requests that are more consistent with the style of the existing code get merged faster. Furthermore, codebases follow common programming patterns or system design patterns that are implicitly specified by the way the current code is written. However, when using code LLMs, these constraints are often inferred incorrectly \citep{wang2024beyond}.
% \naman{good point but would nitpick that models probably learn "style" quite well.. abstractions i think less so?}

% \textit{Example: Serializer-Deserializer pattern for objects:} Consider the issue \href{https://github.com/astropy/astropy/issues/14181}{astropy-\#14181} from the \texttt{astropy} Python library. The issue requests support for a new input file format (reStructuredText) to load astronomical data into the codebase more flexibly. While the issue does not mention it explicitly, as per common practices, developers implement read (deserializer) and write (serializer) operations when implementing support for a new file format. This ensures data can flow bidirectionally between the file format and the application's internal data structures. However, models evaluated on this issue, as part of the SWEBench benchmark, only implemented the \texttt{read} method.

% As discussed in Section \ref{sec:c-hai-colab}, today, unless the specification is overly vague, code LLMs prioritize a functionally correct code snippet without clarifying ambiguities with the user.
% We envision a future where the LLM can proactively surface these decision points and defer back to the user for their resolution.

% \textbf{Tests as specifications}: Another approach to specify software behavior is through tests. These range from input-output examples and assertions to property-based tests. While tests are more concrete than natural language, they can also be underspecified. In practice, test suites are often incomplete and fail to fully capture the desired behavior, especially corner cases. This can result in user misalignment, where AI code passes tests but does not fully satisfy the intended functionality, thus fooling the human. 
% Users may believe tests adequately specify the intended behavior with seemingly obvious examples, yet underspecification can still lead to unexpected behavior that passes these tests. 

\textbf{Formal specifications}: Autoformalization can mitigate underspecification by translating user intent into formal specifications. Ideally, a formal language could completely and unambiguously capture the desired software behavior. However, autoformalization may also suffer from the misalignment problem by misinterpreting user intent, resulting in a formal specification that does not accurately reflect the user's actual requirements. Therefore, users must understand and carefully verify the formalized specification generated by autoformalization tools. 
% If the autoformalization is inaccurate, it introduces another source of potential misalignment, even with the use of formal specifications.


% \todo{Cite: Trade-offs in test suites reduction \citep{shi2014balancing}}

\vspace{-0.2cm}

\subsection{Long-Horizon Code Planning}

\vspace{-0.1cm}

When working on large projects, engineers and tech leads often make complex decisions about how to design and structure the code to best support the various functionalities that will eventually be needed. To build a long-lasting software system, an engineer must know the potential paths that the system's evolution might take. This requires domain expertise and experience in how different code structures require different forms of extension. 
% We remain sceptical about the ability of language models to perform this level of sophisticated planning.

\textbf{Designing good abstractions}:
One instance of long-horizon code planning is choosing the right abstractions from the onset. An API designed with good abstractions will allow new features to be implemented seamlessly with minimal user overhead, while an API designed with poor abstractions may lead to excessive code duplication, refactoring, or even debugging. One example of this is \textit{Library learning}: designing APIs and libraries with useful abstractions often leads to more code reuse and more intuitive interfaces. The challenge of library learning is to derive a library of useful abstractions from a corpus of programs by abstracting out common, reusable computations \cite{ellis2021dreamcoder,10.5555/3692070.3693967}. While the traditional library learning literature has focused primarily on code reuse, a truly effective library must also prioritize ease of use and maintainability, as well as be robust and adaptable to future extensions. %within the software system.
Another challenge is \textit{data representation}, as the choice between data structures leads to a variety of trade-offs when it comes to performance aspects such as memory usage and processing speed. For example, database engineers need to decide between various data models, storage formats, and indexing methods to balance performance. 

% \textit{Example: Database Design for Web Applications}: Database engineers strive to design their databases in a way that minimizes the amount of memory used and maximizes the performance (speed) of their queries. To achieve this goal, the databases community has spent considerable efforts optimizing both the high-level representations of their data and the underlying data structures used \citep{kraska2018case, hawkins2011data}. For this example, let's consider the concrete task of designing a database schema for a restaurant owner to manage various aspects of their business: keeping track of customers, managing a rewards program, maintaining the restaurant's inventory of ingredients, etc. There are a few important design decisions in this task. One of these is the design of the schema: including tables such as reservations and customers is a relatively straightforward decision to make. But should we include a table for customer reviews, or should we just include a rating and review field for each customer? Another important and potentially tricky data representation design decision is choosing which database indexes to include. On one hand, including indexes makes queries significantly faster. On the other hand, indexes take up additional memory and must be updated any time the data is changed. Therefore, in order to design the right set of abstractions, we must be able to estimate the query patterns of the restaurant owners. 
% \naman{the spirit of example is nice but just maybe a bit toyish?}

% \textbf{Code Debugging}:
% \naman{
% Bug-fixing can be particularly challenging since the bugs may present themselves in considerably higher-order effects of the underlying issue.
% For example, \todo{find a swebench or github issue amplifying this aspect}.
% Solving bugs from such bug reports requires reasoning and codebase understanding to localize the cause of such underlying failure. \todo{this should be well-studied in SE research, cite from there?}
% Next, fixing the issue requires rethinking the implementation to fix the issue. 
% These implementation changes should be carefully planned, attempting to come up with \textit{most general} version of the changes while also keeping in mind other parts of the codebase that are affected.
% }


% \textbf{Abstractions in Code Refactoring}: Being able to perfectly predict the evolution of a software project is nearly impossible, and even highly experienced senior engineers with domain expertise fail to design all the abstractions correctly on the first try. Therefore, software often has to undergo major global refactorings. Fundamentally, refactoring is a process of mapping components in the old abstraction to those in the new one in a way that maintains semantic equivalence. This can be difficult for language models because these mappings can be relatively complex: large sections of code can be consolidated or split, the relationships between components may be redefined, and successfully maintaining the semantics properly requires precision. These challenges are reminiscent of those faced by the programming languages community when designing good refactoring tools.
% \naman{hmm i think refactoring is "too concrete" to talk about planning/abstractions -- the abstract bit in refactoring is making design decisions keeping long-term in mind..
% the content here on maintaining mapping seems lower level changes
% }
% First, there is a specification issue, as writing a full specification for a refactoring can be extremely complex, if not impossible. Second, there is a long-context issue, as refactorings can involve 

% In a sense, refactoring is about swapping out one set of abstractions for a new set of abstractions that is easier to use or can better support new use cases. Refactoring can be difficult for LLMs for several reasons. 

%  Finally, refactoring requires a deep understanding of both the original set of abstractions and the new ones. 
 
% \todo{refactoring. highly experienced senior engineers with domain expertise fail to get the abstractions right on the first try. a lot of things have to go through major global refactorings. discuss some of the challenges that are hard to achieve through pure LLM. 1) lots of LOC, 10k+, context window; 2) describing new structure and how it relates to the old structure. this is one of the challenges of PL tools too. 3) hard to express refactoring with a simple specification. MSR agentic refactoring work.  }

% \todo{Other points about long-horizon code planning?}

% \naman{\textbf{Tradeoffs}: I think talking about tradeoffs should be useful here.. rich space with long-term feedback..
% also acknowledging humans don't get it right in first attempt and major software do refactor
% /VLLM just did a major refactoring for v1 release actually!!}
\vspace{-0.2cm}

\subsection{Large Scope and Long Contexts} \label{sec:c-long-context-retrieval}

\vspace{-0.1cm}

The tasks in Sec. \ref{sec:tasks-milestones} become significantly more difficult at the repository-level scope, as engineering with large codebases is a multi-step task. For code generation, there are naturally many decision points where specifications can be vague. Because these decisions compound, it is tricky to generate the right code that the user wants. In code refactoring, modifications will touch multiple parts of the codebase, and it can be tricky to keep the repository consistent. In both code debugging and code navigation, as repositories get bigger, localization becomes more difficult. In code debugging, functions can be large and bugs can be nested deeply within stacks of function calls. In code navigation, because there are so many functions interacting in various ways, it can be difficult to know where each piece of functionality is implemented and how the code is put together.

\textbf{The limits of long-context models}: Software engineering often requires dealing with very large codebases--for example, Google has repositories with over a billion lines of code \citep{potvin2016google}. 
Dealing with large context lengths has long been an outstanding challenge for LLMs, and progress has been rapid with Gemini providing over 1M tokens of context. 
In coding, Magic has even trained models with 100M token context.
% Long context provides an alternative but complementary approach for dealing with large repositories. 
While these models show good performance in toy benchmarks such as needle retrieval, we have not seen evidence of how this performance translates to large real-world code-bases.


% Examples of benchmarks measuring long-context coding capabilities are SWE-Bench \citep{jimenez2024swebench} and Long Code Arena \citep{bogomolov2024long}. These long-context problems are often solved by ...

% \todo{Summarize how long-context coding is dealt with in the literature.}

\textbf{The limits of retrieval-augmented generation (RAG)}: Retrieval-based algorithms have been the predominant way to deal with long-context coding issues. First, the retriever retrieves relevant functions. Then, the generator leverages the retrieval to improve generation. While RAG has proven effective in many NLP tasks such as question answering \citep{gao2023retrieval, lewis2020retrieval}, the code domain provides new challenges for these methods. In the code domain, for the \textit{retrieval step}, it is often unclear what the correct item to retrieve is. In addition, code embeddings generally group code together via syntactic similarity \citep{ma2024unveiling, utpala2023language}, which can make it hard to retrieve crucial snippets that are semantically relevant but syntactically unrelated. For the \textit{generation step}, in NLP tasks it often is a straightforward application of the retrieved information. However, in code, writing a new function requires more than copying and pasting retrieved code snippets. Rather, the relevant snippets must first be analyzed, an algorithm using these code snippets must be developed, and be pieced together coherently in a precise manner.

% \todo{Example of bad generation with good retrieval}

% \naman{\textbf{The limits of long-context models}: 
% Context length for LLMs have been increasing rapidly with Gemini providing over 1 million token context. 
% In the space of coding models, we have witnessed developments of massive context lengths, example over 100 million tokens ~\citep{magic-dev}.
% Long context provides an alternative but complementary approach for dealing with large repositories. 
% These models have been achieving good performance in toy benchmarks such as needle retrieval, however, it is unclear how their performance would scale with real-world large context use cases.
% }

\vspace{-0.25cm}

\subsection{Global Understanding of Codebases} \label{sec:c-global-understanding}

\vspace{-0.1cm}


% \naman{why is this not a part of abstractions? -- comment based on subsec title}

After working with a codebase for a while, programmers have a global understanding and mental model of the codebase. This includes overall codebase structure, different algorithms, stylistic aspects, data representation, program invariants, package versions, test coverage, and the interplay between different functions. This holistic understanding is necessary for performing many tasks.

% A global and holistic understanding of a codebase is important for performing almost all code-related tasks. For example, let's say an engineer is asked to improve the runtime performance of a query. To do so, they must first understand the codebase's structure well enough to know where all the pieces of the algorithm are implemented. Then, they need to understand the algorithm and implementation in detail. This includes both the high-level algorithm (including its time complexity) and the low-level implementation details to identify both algorithmic and implementation bottlenecks. Finally, after coming up with a solution, they must then return to their understanding of the global code structure so that they can integrate their new algorithm without introducing new bugs. 

% \textbf{Understanding invariants for code refactoring}: In order to refactor a large codebase, it is essential to have a complete understanding of how different parts of the codebase interact. These interactions are governed by a set of rules known as program invariants, properties about the code that are assumed to hold. While some of these invariants can be maintained and guaranteed explicitly with assertions, many of these invariants are implicit and may only be obvious to developers working on the codebase itself. When performing code refactoring, developers often have a mental model of these invariants and know what needs to be changed so that these invariants continue to hold. LLMs, on the other hand, lack such a nuanced understanding of programs and may not be able to infer or maintain complex invariants.

% \todo{Example of LLM not being able to understand an invariant.}

% \textbf{Understanding program execution for code debugging}: Debugging code is another task that relies on having a complete global understanding of a codebase. In order to debug programs productively, we need a good internal model of the intermediate execution states of a program. Step-by-step debuggers are generally helpful because the programmer has a mental model of properties that should be true at certain program states, and the debugger can surface whether those properties are indeed true. 

LLMs struggle at global understanding of codebases for several reasons. First, the way that code is pieced together can be relatively intricate, and understanding all these complex relationships can be difficult. Second, code can have units with high algorithmic complexity with custom algorithms that may never have appeared anywhere in the training data. Finally, because a disproportionately large number of LLM training tokens are spent on code generation rather than other coding tasks, models may lack a holistic awareness and world model of code. Generalizing across coding tasks may not be as simple as training: \citet{gu2024cruxeval} found that fine-tuning coding models on additional problems and solutions led to significant improvements on code generation but not code execution. Here, additional training on coding data did not transfer to improving code understanding and execution.
% Although there have been efforts training with execution states to improve general code understanding \citep{ding2024traced}, LLMs still have a poor global model of code.

% \todo{Example of LLM not being to reason about the execution of a piece of code}

% \todo{\url{https://www.expresscomputer.in/news/quantifying-github-copilots-impact-on-code-quality-ai/104480/} cite and also r2e discussion section}

\vspace{-0.25cm}

\subsection{Low-Resource Languages and Custom Libraries} \label{sec:c-low-resource}

\vspace{-0.2cm}

% \todo{Personalization and proprietary code}

As we adapt code LLMs to individual codebases, generating correct code in out of distribution (OOD) scenarios becomes crucial. Much of software development in business contexts revolves around proprietary codebases, a distribution shift from the open-source code that dominates LLM training data  \citep{ahmed2024studying}. Examples include domain-specific languages (DSLs), custom internal libraries, low-resource APIs, and company-specific coding styles/conventions.
% For example, \citet{ahmed2024studying} find a distribution shift between open-source and proprietary C\#/C++ code at Microsoft.

\textbf{Syntactic failures and poor semantic understanding}: Models often hallucinate constructs from higher resource languages when working in low-resource languages. \citet{blinn2024statically} found that when writing Hazel, LLMs often borrowed syntax and library functions from OCaml and Elm, higher-resource languages. When writing Triton, models often use syntactically incorrect constructs such as array indexing.
In addition, models have less exposure to the various language constructs. Therefore, they have a weaker semantic understanding of the language. Many studies reveal that code LLMs perform poorly when asked to write code in low-resource languages. Due to the lack of training data in these OOD domains, models may struggle to write common primitives or piece together functionality coherently. On HumanEval, Qwen 2.5 Coder Instruct (32B) \citep{hui2024qwen2} has an accuracy of 83\% in Python but only 27\% in D.\footnote{As reported by the \href{https://huggingface.co/spaces/bigcode/bigcode-models-leaderboard}{BigCode Models Leaderboard} on the MultiPL-E benchmark \citep{cassano2023multipl}}

% \todo{Example of LLM not understanding low-resource language semantics}

% \textbf{Domain-Specific Languages (DSLs)}: Similar to low-resource programming languages, generating programs in a DSL is also challenging for LLMs due to the out-of-distribution nature of this task. In Section 6.5 of \citet{qiu2024tenspiler}, the authors find that Claude Opus struggles to generate code in three DSLs, making mistakes such as hallucinating libraries that do not exist, including the wrong imports, and using APIs improperly. For a more comprehensive survey of low-resource languages and DSLs, including current progress in these directions, we refer the reader to \citet{joel2024survey}. 

\textbf{Library usage failures}: In OOD scenarios, LLMs lack awareness of the libraries and functions available for use. In new codebases using custom libraries, many functions appear only a few times, providing limited training data for AI models to learn their usage. This scarcity can lead to overfitting, where models fail to recognize an effective use-case of these functions. Models also frequently hallucinate non-existent functions based on patterns that it infers.

% In Listing \ref{lst:lean-hallucination}, the LLM makes up a theorem named \texttt{infinitelyManyPrimes\_arithmeticSequence}. While the model understands that it should use Dirichlet's theorem, the actual theorem is named \texttt{forall\_exists\_prime\_gt\_and\_eq\_mod}. This theorem occurs very few times in all of the currently available Lean code on GitHub, so it is likely that LLMs did not train on it many times.

% \begin{lstlisting}[caption={LLM-generated Lean code snippet that incorrectly uses indexing notation},label={lst:lean-hallucination}, captionpos=t, language=lean, breaklines=true, style=mystyle]
% theorem infinitely_many_primes_of_form_6k_plus_1 :
%   ∀ n : ℕ, ∃ p : ℕ, p > n ∧ Prime p ∧ ∃ k : ℕ, p = 6 * k + 1 := by
%   intro n
%   -- Use Dirichlet's theorem for arithmetic progressions
%   !*have h := Nat.infinitelyManyPrimes_arithmeticSequence 6 1
% \end{lstlisting}

% OOD \citep{zhou2023devil}

% high resource intermediate language \citep{mora2024synthetic}


% \begin{itemize}
%     \item Low-resource languages.
%     \item Different programming paradigms, like imperative vs declarative.
%     \item Even for high-resource languages, can models handle OOD scenarios well?
%     \item Low-resource library usage.
%     \item Codebase-specific style guidelines.
%     \item There are no idiomatic evaluations.
% \end{itemize}

% \todo{relationship between OOD/custom libraries and code reuse case. code reuse is where you have a set of building blocks and need to combine them, the difficulty is that the amount of blocks is very large.}

\vspace{-0.2cm}

\subsection{Library and API Version Updates} \label{sec:c-library-updates}

\vspace{-0.1cm}

Continual learning, the idea of training an AI system to take in new information continually, has been a long-standing challenge in AI and NLP  \citep{wu2024continual, wang2024comprehensive}. In software engineering, codebases are continuously changing as new features are supported and awkward design patterns are reworked. While backwards compatibility is often prioritized in software design, it inevitably becomes broken as codebases evolve further. Therefore, programming libraries have version releases, each release supporting and deprecating features in the last version. 

% Benchmarks such as CodeUpdateArena \citep{liu2024codeupdatearena} and GitChameleon \citep{islah2024gitchameleon} examine this issue at the function and file level, finding that language models already struggle to adapt to these changes even with this limited scope. In theorem proving, \citet{kumarappan2024leanagent} aim to mitigate this by developing a lifelong learning framework that continuously learns and uses new theorems. In real-world engineering, however, the challenge of library and API versioning generally spans across an entire repository, as everything must be kept consistent. 
% Currently, there are no benchmarks or techniques that successfully deal with this challenge at such a large scale. 

% \textbf{Version identification}: It is difficult for LLMs to even know which version of each library is being used in a codebase. This information can be hidden deeply within a codebase (e.g., in a configuration file). Even worse, most code is compatible across multiple versions, certain features will break in specific versions. The model will need a deep understanding of the nuances between different versions to infer which version is being used. 

\textbf{Continous adaptation}: Language models are not as dynamic as codebases, and new features and APIs released yesterday will not be in the training data of today's code LMs. Code written with new APIs and new versions of existing libraries are low-resource. Even when there are new and simplified ways to write features, LLMs may rely on older, more cumbersome approaches because those are the ones seen more frequently in the model's training data. 

\textbf{Version-specific constructs}:
For fast-changing libraries that are not backward compatible, it can be difficult for LLMs to implicitly keep track of which constructs are associated with each version. This makes consistently using constructs from the right version difficult. Therefore, LLMs may write code that mixes and matches API constructs from different versions of the same library. 

% \todo{Perhaps add probing tasks to see if LLMs explicitly know which constructs are in which library}

% \todo{Example: Jane Street OCaml vs. general OCaml}

% \textit{Example. Lean 3 vs. Lean 4}: Lean \citep{de2015lean} is a programming language that allows users to write formal proofs of mathematical theorems. In 2017, using Lean 3, enthusiasts implemented a library for mathematics called \texttt{mathlib}, with over half a million lines of code. Because Lean 3 had many shortcomings, Lean 4 \citep{moura2021lean} was initiated at the beginning of 2021 to address many of these issues. There was a massive undertaking to port all of the \texttt{mathlib} code over to Lean 4, and only in September 2023 was there a stable release of Lean 4, the version of Lean that is predominantly used today. The two versions are generally incompatible. We hypothesize that, due to the recency of Lean 4, most language models have been trained on much more Lean 3 code compared to Lean 4 code. This can be seen when asking language models to generate code in Lean 4, as it sometimes generates code with Lean 3 coding conventions or uses theorems and lemmas from Lean 3 that are deprecated in Lean 4. 

% \textit{Example. Typing Hints}: While Python 3.5 required importing types from the typing module, Python 3.9's PEP 585 enabled direct use of built-in types for generics (e.g., \texttt{\small list[int]} vs \texttt{\small typing.List[int]}). However, language models tend to default to the older \texttt{typing} module syntax.

% \textit{Example. Dynamo}:  \todo{Armando's example: Copilot is bad at dynamo tasks because there are a bunch of versions of the API so mixes things up.}

\vspace{-0.2cm}

\subsection{High Algorithmic Complexity: OOD Domains} \label{sec:c-ood-domains}

\vspace{-0.2cm}

Some programming tasks are challenging for even the best human programmers, requiring approaches with a very high algorithmic complexity. Examples of tasks that fall into this category include superoptimizing programs, discovering attacks for purportedly secure code, writing performant compilers, optimizing GPU kernels \citep{ouyang2024kernelbench}, and writing very error-prone and very technical code. 
%
% \todo{Flesh out the difficulties more}
% \todo{Add more examples of superhuman / OOD domains}
% \textit{Example. Synthesis of Sorting Kernels}: An example of an out-of-distribution domain is synthesizing fast assembly code for sorting kernels. In 2023, AlphaDev \citep{mankowitz2023faster} used reinforcement learning to find a SoTA kernel for sorting length 3-5 arrays. While this appeared to be a superhuman performance, shortly after, \cite{neri2023shorter} hand-wrote a kernel shorter and faster than the one found by AlphaDev. Later, \citep{ullrich2024synthesis} developed an algorithm based on enumeration and intelligent heuristic-based sampling that beat both of these. In addition, the algorithm ran faster than AlphaDev by two orders of magnitude. In this case, while AI was able to achieve an impressive performance, humans were able to discover better algorithms.
%
% \textbf{Limits of symbolic techniques}: When it comes to applying symbolic techniques to these tasks, there are a few limiting factors that make them difficult to tackle. First, for synthesis-style tasks, the search space can be very large. Deductive and rewrite-based synthesis techniques are unable to explore a majority of the search space. Second, verifiers can be limited in power, such as when dealing with properties in concurrency or weak memory models. Third, many domains lack clean models to reason about properties, such as dealing with memory bandwidth in GPU kernels. 
%
Because they are hard for humans, these tasks are very rarely in the training data of today's language models. They have unique, domain-specific, challenges that making generalizing from existing data difficult. For these problems, language models rely heavily on feedback-driven search algorithms \citep{mankowitz2023faster}, and it can be difficult to navigate the search space effectively. In addition, many of these tasks lack feedback mechanisms, which is crucial for AI to pick up learning signals. When designing a complex algorithm or data structure, it is often hard to know if you are on the right track until you get to the correct result. When writing code for a large multithreaded operation, it may be hard to know if the algorithm has concurrency issues until all the parts are fully fleshed out. Without feedback, incremental improvement is nearly impossible.


% \begin{itemize}
%     \item New systems/languages (particularly excited for supporting verification where from very limited understanding bottleneck is good general systems since domain experts write one of things (or proprietary)).
% \end{itemize}